% chapters/ch21_separation_principle.tex
% Chapter 21: NRMO-Engine Separation Principle — Formal Definition
% Source: NRMO_v5_updated.pdf, Chapter 21 (lines 9324-9433 of v5_layout.txt)

\chapter{NRMO-Engine Separation Principle ---
Formal Definition}
\label{ch:separation-principle}

\nrmobox{Document status}{
\begin{tabular}{ll}
\textbf{Status}: & Structural Requirement \\
\textbf{Precondition}: & Strong Engine Architecture
(Chapter~\ref{ch:strong-engine})
\end{tabular}\\[0.3em]
This chapter formalises the separation between NRMO (constitution)
and Engine (propulsion) as an invariant.
}

% =====================================================================
\section{Prohibited Patterns (Failed in All World Lines)}
\label{sec:sep-prohibited}

\subsection{Pattern 1: Engine Containment}

\begin{lstlisting}[basicstyle=\ttfamily\small,frame=single]
Engine subset-of NRMO

Result:
  - Non-continuous jump: impossible
  - Cannot cross threshold
  - Loss of initiative

Proof: Contradiction with P(Ruin) <= epsilon constraint
\end{lstlisting}

\subsection{Pattern 2: Intervention in Purpose / Strategy}

\begin{lstlisting}[basicstyle=\ttfamily\small,frame=single]
NRMO -> Engine.purpose
NRMO -> Engine.strategy
NRMO -> Engine.success_definition

Result:
  - Constitution becomes ruler
  - Disappearance of decision subject
  - Outsourcing of responsibility
\end{lstlisting}

\subsection{Pattern 3: Exploration Suppression}

\begin{lstlisting}[basicstyle=\ttfamily\small,frame=single]
Prohibit Engine.exploration under "safety" pretext

Result:
  - Religionization (absolutizing NRMO)
  - Stagnation (disappearance of exploration)
  - Self-exculpation
\end{lstlisting}

% =====================================================================
\section{Separation Principle (Formal Definition)}
\label{sec:sep-formal}

\begin{lstlisting}[basicstyle=\ttfamily\scriptsize,frame=single]
Separation_Principle = {
  "placement": {
    "Engine not-subset-of NRMO":  True,
    "Engine.location":            "Outside(NRMO)"
  },
  "authority": {
    "NRMO.authority":     "veto_only",
    "NRMO.authority_NOT": [
      "propulsion", "optimization",
      "success_definition",
      "purpose_setting",
      "strategy_decision"
    ]
  },
  "relationship": {
    "NRMO -> Engine":   "veto possible (one-way)",
    "NRMO !-> Engine":  "purpose/strategy intervention prohibited",
    "Engine -> NRMO":   "cannot ignore"
  },
  "health": {
    "health":       "inversely proportional to veto_frequency",
    "pathological": [
      "high-frequency veto",
      "exploration suppression",
      "becoming ruler"
    ]
  }
}
\end{lstlisting}

% =====================================================================
\section{Verification Method at Implementation}
\label{sec:sep-verification}

\subsection{Self-Check Questions}

\begin{enumerate}
\item Is Engine outside NRMO?
\item Is NRMO not instructing ``what should be done''?
\item Is NRMO not defining ``how to succeed''?
\item Is NRMO exercising veto authority only?
\item Is intervention frequency decreasing over time?
\end{enumerate}

\noindent
All ``Yes'' = healthy. Even one ``No'' = suspected structural
violation.

% =====================================================================
\section*{Connections to Other Parts}

\begin{itemize}
\item This chapter formalises Invariant~\ref{inv:gov-exec-sep}
(governance--execution separation).
\item In the engineering implementation (Part~X), the separation
principle is verified by the audit checks documented in
\texttt{SEPARATION.md} of the v5.2 codebase: \emph{(i)} Engine
modules contain no governance parameters; \emph{(ii)} flow order is
\emph{invent $\to$ governance filter $\to$ engine scoring};
\emph{(iii)} \texttt{omega\_full.py} delegates all ruin checks to
\texttt{core.ruin.is\_ruin\_state()}.
\item In $\Omega$ Full v5.0 (Part~IX), the separation principle
appears as the architectural invariant: \emph{NRMO defines
admissibility; engine searches within it.}
\end{itemize}
